\section{Introduction}
\label{sec:introduction}

The integration of Internet of Things (IoT) technologies and artificial intelligence has initiated a paradigm shift in modern precision agriculture, transforming traditional fields into complex cyber-physical systems \cite{Bayar2025, PourgholamAmiji2024}. Central to this transformation is the optimization of irrigation, a critical requirement for mitigating regional water scarcity while maintaining agricultural productivity \cite{Nouri2023, Maniam2026}. While advanced optimal control algorithms like Model Predictive Control (MPC) and Reinforcement Learning (RL) are highly capable of managing these constrained multivariate systems, deploying them directly into physical environments is economically risky and biologically prohibitive. Consequently, the synthesis and training of these intelligent controllers necessitate a high-fidelity virtual environment that accurately captures the complex, non-linear dynamics of the soil-plant-atmosphere continuum.

This work presents the complete physical and mathematical description of an agricultural system designed based on a land in northwestern Iran. The objective is to construct a high-fidelity simulation environment that serves a dual purpose: first, as the virtual plant on which the Model Predictive Controller operates, and second, as the training environment for the Reinforcement Learning agent later developed.

The modeling framework builds upon the agent-based model (ABM) proposed by L\'{o}pez-Jim\'{e}nez et al.~\cite{Lopezjimenez2024}, which couples soil water dynamics with crop growth on a topographically heterogeneous field. However, several critical modifications and extensions are introduced to adapt the model to the specific agronomic conditions of rice cultivation in Gilan Province, Iran, and to correct certain mass-balance inconsistencies identified in the original formulation.

The report is organized as follows. Section~\ref{sec:study_site} describes the physical field site and the digital elevation model used to define the directed flow graph and agent grid. Section~\ref{sec:climate} details the 25-year climate dataset, the reference evapotranspiration modeling, and the contrasting meteorological scenarios. Section~\ref{sec:abm_equations} presents the core discrete-time ABM equations governing the decoupled surface and subsurface water balances, thermal accumulation, and crop growth. Section~\ref{sec:corrections} formally documents the necessary mass-conservation corrections and agronomic extensions applied to the baseline framework. Section~\ref{sec:soil_crop_params} parameterizes the environment for Hashemi rice and tobacco. Section~\ref{sec:water_budget} derives the agronomically grounded seasonal water budgets. Finally, Section~\ref{sec:validation} presents the ABM validation results, including cross-validation against satellite soil wetness observations. Section~\ref{sec:conclusion} concludes the work.

% ========================= SECTION 1 ================================= %
\newpage
\section{Study Site and Topological Graph}
\label{sec:study_site}

The study site is located on the slopes of the Talish Mountains in Gilan Province in northwest Iran, at coordinates $38.298^{\circ}$N, $48.847^{\circ}$E. Situated along the southwestern coast of the Caspian Sea, Gilan is Iran's primary rice-producing province, accounting for over 238,000 hectares of paddy fields \cite{surfiran2025}. The province experiences a humid subtropical climate with hot summers and relatively mild winters, although seasonal precipitation can vary dramatically between years. Despite its comparatively favorable water endowment relative to Iran's central plateau, Gilan faces mounting pressures from declining aquifer levels and increasing competition between agricultural, urban, and industrial water users~\cite{Nouri2023,Mesgaran2018}. These conditions make this land a representative case study for evaluating constrained irrigation control strategies.

\begin{figure}[H]
    \centering
    \includegraphics[width=0.75\textwidth]{figures/land.png}
    \caption{The satellite image of the study site captured using Google Earth}
    \label{fig:land}
\end{figure}

\subsection{Digital Elevation Model}
\label{sec:field}

The selected field covers approximately 6 hectares and is discretized into a $10 \times 13$ rectangular grid, yielding $N = 130$ crop-soil agents. This resolution represents a practical balance between spatial detail and computational tractability for the simulation environment: each agent corresponds to a plot of roughly 460~m$^2$, which is sufficiently fine to capture the field's topographical heterogeneity while remaining computationally feasible for the MPC optimization solved at every daily time step.

The topographical data were obtained from the USGS Earth Explorer digital elevation database which provides Shuttle Radar Topography Mission (SRTM) data at 30-meter resolution. The raw elevation raster was processed in QGIS 3034 to extract the study region, reproject it to
UTM Zone 38N, and resample it to the $10 \times 13$ agent grid. The resulting digital elevation model (DEM) is stored as a GeoTIFF file and loaded programmatically by the simulation environment.

\begin{figure}[H]
    \centering
    \includegraphics[width=0.75\textwidth]{figures/elevation_map.png}
    \caption{Three-dimensional surface plot of the 130-agent grid. Elevation ranges from 70~m to 181~m above sea level, forming a bowl-shaped topology with water convergence toward the southeastern corner.}
    \label{fig:topo3d}
\end{figure}

The elevation field exhibits a distinctive bowl-shaped topography, with elevations ranging from 70~m above sea level in the southeastern corner to to 181~m in the northwestern and southwestern ridges. This 111~m elevation gradient creates significant lateral water redistribution through surface and subsurface runoff: agents on higher ground shed excess water to their downhill neighbors, while agents in the valley bottom accumulate inflows from multiple directions. This topographical heterogeneity is a defining feature of the study site and the principal motivation for adopting an agent-based rather than a lumped-parameter modeling approach.

\subsection{Directed Graph Construction}
\label{sec:directed_graph}

The topographical information is encoded in a directed graph that governs lateral water exchange between agents. Each agent occupies a cell in the $10 \times 13$ grid and is assigned a normalized elevation $\gamma^{(n)} \in [0, 1]$, computed by min--max normalization of the raw DEM values:
\begin{equation}
    \gamma^{(n)} = \frac{z^{(n)} - z_{\min}}{z_{\max} - z_{\min}}
    \label{eq:gamma}
\end{equation}

where $z^{(n)}$ is the raw elevation of agent $n$. An agent at position $(r, c)$ can exchange water with up to eight neighbors in its Moore neighborhood (the eight surrounding cells). A directed edge from agent $n$ to agent $m$ is created if and only if $\gamma^{(m)} < \gamma^{(n)}$, so that water flows strictly downhill. The number of lower-elevation neighbors for each agent is denoted $N_r^{(n)}$.

\begin{figure}[H]
    \centering
    \includegraphics[width=0.75\textwidth]{figures/directed_graphs.png}
    \caption{Water flow directions between agents, overlaid on the normalized elevation map. Arrows indicate the direction of lateral surface runoff, governed by the directed graph.}
    \label{fig:flow_arrows}
\end{figure}

Figure~\ref{fig:flow_arrows} visualizes the resulting directed graph. The arrow density reflects the terrain gradient: steep areas near the ridges exhibit dense, nearly parallel flow lines, while the flat valley floor shows sparser, more diffuse connections. Agents in the southeastern corner with $N_r = 0$ function as hydrological sinks; they have no lower neighbors and therefore retain all accumulated water. These sink agents are particularly susceptible to waterlogging stress, a phenomenon that the ABM captures through the waterlogging stress function $h_6$ described in Section~\ref{sec:stress_growth}.

\begin{figure}[H]
    \centering
    \includegraphics[width=0.7\textwidth]{figures/sinks.png}
    \caption{Spatial distribution of the neighbor count $N_r$ for each agent. Red-bordered cells indicate sink agents ($N_r = 0$) that accumulate runoff from uphill neighbors.}
    \label{fig:sink_agents}
\end{figure}

Three sink agents were identified, two located in the southeastern corner of the field at the lowest elevations (Figure~\ref{fig:sink_agents}). This spatial concentration of sinks creates a natural drainage basin structure that any irrigation controller must account for: over-irrigating uphill agents will cause water to pool at these sinks, simultaneously wasting budget and inducing waterlogging stress in the valley.

% ========================= SECTION 2 ================================= %

\section{Climate Data and Scenarios}
\label{sec:climate}

Daily meteorological data were obtained from the NASA POWER (Prediction of Worldwide Energy Resource) database, which provides satellite-derived and reanalysis-based climate records at a spatial resolution of $0.5 \times 0.625$ (approximately $50 \times 60$~km). The dataset was requested for 25-year period of 2000--2025 at the field coordinates ($38.30^{\circ}$N, $48.85^{\circ}$E). Thirteen variables were extracted, including 

\begin{itemize}
    \item mean, maximum, and minimum temperatures at 2~m (T2M, T2M\_MAX, T2M\_MIN, ${\circ}$C), 
    \item dew/frost point at 2~m (T2MDEW, ${\circ}$C), 
    \item corrected daily precipitation (PRECTOTCORR, mm/day), 
    \item relative humidity (RH2M, \%), 
    \item surface shortwave downward irradiance (ALLSKY\_SFC\_SW\_DWN, MJ/m$^2$/day), 
    \item wind speed at 2~m (WS2M, m/s), 
    \item surface pressure (PS, kPa), and 
    \item root and surface soil wetness (GWETROOT, GWETTOP, dimensionless, 0–1). 
\end{itemize}

The last two variables were used for validation purposes rather than as model inputs.

The dataset was filtered to cover the agricultural season from April to October for each year, yielding 5,564 daily records with no missing values. During preprocessing, one extreme rainfall anomaly (375~mm recorded on April 7, 2022) was identified and replaced with the calendar-day average of 1.2~mm to prevent anomalous skewing of the soil water balance. Moreover, based on the available data, two reference evapotranspiration values were calculated and discussed in the following subsection.

\subsection{Reference Evapotranspiration}
\label{sec:et0}

Reference evapotranspiration (ET$_0$) is the primary driver of crop water demand. Two estimation methods were compared to determine the most accurate representation for the study site: the temperature-based Hargreaves equation and the FAO Penman-Monteith equation \cite{allen1998}. 

\textbf{Hargreaves equation}: 
\begin{equation}
    \text{ET}_0^{H} = 0.0023 \cdot (T_{\text{mean}} + 17.8) \cdot (T_{\text{max}} - T_{\text{min}})^{0.5} \cdot R_a
\end{equation}
where $T_{\text{mean}}$ is the mean daily temperature ($^{\circ}$C), $T_{\max}$ and $T_{\min}$ are the daily maximum and estimated minimum temperatures ($^{\circ}$C), and $R_a$ is the extraterrestrial radiation (expressed in equivalent evaporation, mm/day).

\textbf{FAO Penman-Monteith}:
\begin{equation}
    \text{ET}_0^{PM} = \frac{0.408 \,\Delta\,(R_n - G) + \gamma \,\frac{900}{T+273}\, u_2\, (e_s - e_a)}{\Delta + \gamma\,(1 + 0.34\, u_2)}
\end{equation}

where $\Delta$ is the slope of the vapor pressure curve (kPa/$^{\circ}$C), $R_n$ is the net radiation at the crop surface (MJ/m$^2$/day), $G$ is the soil heat flux density (MJ/m$^2$/day), $\gamma$ is the psychrometric constant (kPa/$^{\circ}$C), $T$ is the mean daily air temperature at 2~m height ($^{\circ}$C), $u_2$ is the wind speed at 2~m height (m/s), $e_s$ is the saturation vapor pressure (kPa), $e_a$ is the actual vapor pressure (kPa), $(e_s - e_a)$ is the saturation vapor pressure deficit (kPa).

Over the 25-year rice season, the Hargreaves method yielded an average ET$_0$ of 4.59~mm/day, compared to 5.02~mm/day for the Penman-Monteith method ($r = 0.940$). Although the models are strongly correlated, the Hargreaves equation systematically underestimates evaporative demand by approximately 8.6\%. This underestimation occurs because the Hargreaves formulation relies on the diurnal temperature range ($T_{\max} - T_{\min}$) as a proxy for atmospheric demand. In Gilan's coastal microclimate, the moderated temperature differential ($\Delta T \approx 9.4^{\circ}$C) artificially attenuates the Hargreaves estimate. Therefore, the Penman-Monteith method, which explicitly incorporates solar radiation, relative humidity, and aerodynamic conductances, 

is selected as the standard to ensure high-fidelity calculation of the evaporative disturbances in the control environment.

\begin{figure}[H]
\centering
\includegraphics[width=\textwidth]{figures/fig_et0_new.png}
\caption{Comparison of Hargreaves and Penman-Monteith ET$_0$. Left: seasonal pattern. Right: scatter plot with regression.}
\label{fig:et0}
\end{figure}

\subsection{25-Year Climatology and Scenario Selection}

The 25-year climatology reveals mean temperatures ranging from $12^{\circ}$C in early April to $26^{\circ}$C in mid-July. Precipitation exhibits a dry mid-season between June and August (0.5--1.5~mm/day) flanked by wetter months, while solar radiation peaks at 26~MJ/m$^{2}$/day in June. Relative humidity remains consistently between 60\% and 75\% throughout the season, which is a defining feature of Gilan's Caspian climate. Additional climate figures are provided in Appendix ~\ref{apdx:plots}. 

\begin{figure}[H]
    \centering
    \includegraphics[width=\textwidth]{figures/fig_temperature.png}
    \caption{25-year daily average temperature (April--October, 2000--2025).}
    \label{fig:temperature}
\end{figure}

\begin{figure}[H]
    \centering
    \includegraphics[width=\textwidth]{figures/fig_rainfall.png}
    \caption{25-year daily average rainfall (April--October, 2000--2025).}
    \label{fig:rainfall}
\end{figure}

To evaluate the irrigation controllers under varying weather conditions, three representative years were selected from the 25-year record to serve as distinct climate scenarios:


\begin{itemize}
    \item \textbf{Moderate Scenario (2020):} Represents baseline conditions.
    \item \textbf{Dry Scenario (2022):} Evaluates controller performance under severe water constraints.
    \item \textbf{Wet Scenario (2024):} Tests the controller's ability to capitalize on natural rainfall and mitigate waterlogging risks in low-elevation agents.
\end{itemize}

\begin{table}[H]
\centering
\caption{Selected climate scenarios.}
\begin{tabular}{lccc}
\toprule
\textbf{Scenario} & \textbf{Year} & \textbf{Rice (mm)} & \textbf{Tobacco (mm)} \\
\midrule
Dry & 2022 & 39.7 & 54.0 \\
Moderate & 2020 & 42.1 & 73.3 \\
Wet & 2024 & 176.8 & 163.2 \\
\bottomrule
\end{tabular}
\end{table}

\section{Agent-Based Model Equations}
\label{sec:abm_equations}

The crop-soil system is modeled as a set of $N=130$ coupled discrete-time dynamical systems. Each agent $n \in \{1, \dots, N\}$ maintains five state variables that evolve at a daily time step $k$: volumetric soil water content in the root zone ($x_1$, mm), cumulative thermal time ($x_2$, $^{\circ}$C$\cdot$day), cumulative maturation index ($x_3$), aboveground dry biomass ($x_4$, g/m$^2$), and surface ponding depth ($x_5$, mm).

To accurately simulate lateral water redistribution while maintaining computational efficiency for the optimal controller, the grid is solved sequentially at each time step $k$ following a topological sort of the directed elevation graph, executing from the highest elevation agents down to the sinks. Furthermore, horizontal subsurface lateral flow is considered negligible for the 460~m$^2$ agent resolution; thus, agents are coupled exclusively via surface runoff.

\subsection{Surface Hydrology and Runoff}

The total available surface water ($W_{\text{surf}}$) for agent $n$ consists of yesterday's un-infiltrated ponding ($x_5$), daily precipitation ($P$), applied irrigation ($u^n$), and incoming surface runoff routed from all uphill neighbors $m \in \mathcal{U}(n)$, divided equally among the neighbor's downhill receivers ($N_r^m$):
\begin{equation}
    W_{\text{surf}}^n(k) = x_5^n(k) + P(k) + u^n(k) + \sum_{m \in \mathcal{U}(n)} \frac{\varphi_2^m(k)}{N_r^m}
\end{equation}

Surface runoff ($\varphi_2$) is computed using the SCS curve number method. To enforce strict mass conservation and simulate topographical valley flooding, a boundary condition is applied where sink agents ($N_r^n = 0$) cannot generate runoff:
\begin{equation}
    \varphi_2^n(k) = 
    \begin{cases}  
      \frac{(W_{\text{surf}}^n(k) - \theta_3)^2}{W_{\text{surf}}^n(k) + 4\theta_3}, & N_r^n > 0 \text{ and } W_{\text{surf}}^n(k) > \theta_3 \\
      0, & \text{otherwise} 
   \end{cases}
\end{equation}

\subsection{Subsurface Hydrology and Infiltration}

Crop transpiration ($\varphi_1$) is the water extracted by plant roots, bounded by the specific crop evapotranspiration ($\text{ET}_c$), which is derived from the atmospheric reference demand ($\text{ET}_0$) and the crop coefficient ($K_c$):
\begin{equation}
    \text{ET}_c(k) = K_c \cdot \text{ET}_0(k)
\end{equation}
\begin{equation}
    \varphi_1^n(k) = \min\!\Big(\max\!\big(\theta_1(x_1^n(k) - \theta_2\theta_5),\; 0\big),\; \text{ET}_c(k)\Big)
\end{equation}

Water that does not run off the surface infiltrates into the root zone. To prevent non-physical infinite pooling within the soil matrix, infiltration ($I$) is bounded by the physical saturation capacity of the soil ($\theta_{\text{sat}}\theta_5$):
\begin{equation}
    I_{\max}^n(k) = \theta_{\text{sat}}\theta_5 - x_1^n(k) + \varphi_1^n(k)
\end{equation}
\begin{equation}
    I^n(k) = \min\!\big(W_{\text{surf}}^n(k) - \varphi_2^n(k),\; I_{\max}^n(k)\big)
\end{equation}

An intermediate soil moisture state ($x_{1,\text{temp}}$) is computed to evaluate the temporary water holding before daily drainage occurs:
\begin{equation}
    x_{1,\text{temp}}^n(k) = x_1^n(k) + I^n(k) - \varphi_1^n(k)
\end{equation}

If this temporary moisture exceeds the field capacity ($\theta_6\theta_5$), the subsurface excess ($E_{\text{sub}}$) becomes vulnerable to gravity. A fraction of this excess, defined by the daily drainage rate ($\theta_4$), is lost to deep vertical drainage ($\varphi_3$):
\begin{equation}
    E_{\text{sub}}^n(k) = \max\!\big(x_{1,\text{temp}}^n(k) - \theta_6\theta_5,\; 0\big)
\end{equation}
\begin{equation}
    \varphi_3^n(k) = \theta_4 \cdot E_{\text{sub}}^n(k)
\end{equation}

\subsection{Stress Penalties and Crop Growth}
\label{sec:stress_growth}

Thermal time accumulation ($h_1$) drives phenological development:
\begin{equation}
    h_1^n(k) = \max\!\big(T_{\text{mean}}(k) - \theta_7,\; 0\big)
\end{equation}

Heat stress ($h_2$) reduces maturation efficiency, and low-temperature stress ($h_7$) prevents biomass accumulation:
\begin{equation}
    h_2(k) = \text{clip}\!\left( \frac{\theta_{10} - T_{\max}(k)}{\theta_{10} - \theta_9},\; 0,\; 1 \right)
\end{equation}
\begin{equation}
    h_7(k) = 
    \begin{cases} 
      1, & T_{\text{mean}}(k) > \theta_7 \\ 
      0, & T_{\text{mean}}(k) \leq \theta_7 
   \end{cases}
\end{equation}

Hydrological stresses directly impact crop development. Drought stress ($h_3$) penalizes the system when actual transpiration falls below the crop's atmospheric demand ($\text{ET}_c$), while waterlogging stress ($h_6$) penalizes root-zone oxygen deprivation when soil moisture exceeds field capacity:
\begin{equation}
    h_3^n(k) = 1 - \theta_{14} \cdot \max\!\left(1 - \frac{\varphi_1^n(k)}{\text{ET}_c(k)},\; 0\right)
\end{equation}
\begin{equation}
    h_6^n(k) = \text{clip}\!\left(1 - \frac{x_1^n(k) - \theta_6\theta_5}{\theta_6\theta_5},\; 0,\; 1\right)
\end{equation}

The growth function ($g$) dictates the fraction of intercepted solar radiation contributing to biomass, forming a two-branch sigmoid curve scaled by the maximum interception parameter ($\theta_{19}$):
\begin{equation}
g(x_2^n) = \begin{cases}
\theta_{19}/(1 + e^{-0.01(x_2^n - \theta_{20})}) & x_2^n \leq \theta_{18}/2 \\
\theta_{19}/(1 + e^{0.01(x_2^n + x_3^n - \theta_{18})}) & x_2^n > \theta_{18}/2
\end{cases}
\end{equation}


\subsection{State Update Equations}

The coupled dynamics for agent $n$ conclude with the daily state updates. Any surface water that failed to infiltrate remains explicitly conserved in the surface ponding state ($x_5$) for the subsequent day:
\begin{align}
    x_1^n(k+1) &= x_{1,\text{temp}}^n(k) - \varphi_3^n(k) \\
    x_2^n(k+1) &= x_2^n(k) + h_1^n(k) \\
    x_3^n(k+1) &= x_3^n(k) + \theta_{11}\big(1 - h_2(k)\big) + \theta_{12}\big(1 - h_3^n(k)\big) \\
    x_4^n(k+1) &= x_4^n(k) + \theta_{13} \cdot h_3^n(k) \cdot h_6^n(k) \cdot h_7(k) \cdot g(x_2^n(k)) \cdot R_s(k) \\
    x_5^n(k+1) &= W_{\text{surf}}^n(k) - \varphi_2^n(k) - I^n(k)
\end{align}

\section{Model Corrections and Extensions}
\label{sec:corrections}

While the baseline framework~\cite{Lopezjimenez2024} provides a robust topographical architecture, implementation for closed-loop control revealed necessary mathematical corrections to enforce strict mass-balance and to ensure controller cost-functions were properly coupled to agronomic penalties. Five primary modifications were integrated into the environment.

\subsection{Hydrological Decoupling and Mass Conservation}

The original formulation in~\cite{Lopezjimenez2024} computes an outflow term ($\phi_{\text{out}}$) that lumps surface runoff and subsurface excess together. While the runoff portion is properly subtracted from the sender's state equation, the subsurface excess portion is not. Consequently, a sending agent passes its subsurface excess to downhill neighbors but never loses that volume from its own root zone. This generates an unbounded artificial creation of water ("phantom water"), leading to exaggerated flooding in lower elevations during sustained rainfall. 

The proposed decoupled two-layer architecture resolves these flaws. By explicitly bridging the surface layer ($W_{\text{surf}}$) and the root zone ($x_1$) via a strictly bounded infiltration term ($I$), the overlapping fluxes are eliminated. Explicitly tracking the exact volume of water moving between the decoupled layers guarantees strict mass conservation across the entire field.

\subsection{Surface Routing, Sinks, and Ponding Dynamics}
\label{sec:runoff}

In the original configuration, surface runoff ($\varphi_2$) is mathematically aggregated with subsurface flow and routed directly into the soil matrix of adjacent agents. This simplification bypasses surface pooling dynamics, injecting rapid overland flow directly into the root zone of downhill neighbors. 

To more accurately capture the bowl-shaped topography of the field, surface runoff is explicitly decoupled from subsurface flow in this work. Runoff generated by an agent is routed to lower-elevation neighbors where it arrives as surface inflow. This inflow is added to the neighbor's available surface water ($W_{\text{surf}}$) and must undergo infiltration before entering the root zone. 

Crucially, a topological boundary condition was introduced to handle valley flooding. In the baseline model, sink agents ($N_r = 0$) generated runoff that mathematically vanished. By enforcing $\varphi_2 = 0$ for all sinks, and by capturing un-infiltrated water in a new surface ponding state variable ($x_5$), the model realistically simulates sustained, multi-day waterlogging in low-elevation zones.

To justify the extra computational cost of routing water across the elevation map, the cascade model was compared against a simpler single-step baseline where each grid cell acts independently. Both models were tested under the wet year scenario (2024) with three different irrigation levels (0 mm, 5 mm, and 8 mm). While the simpler model produces similar field-average yields, the cascade mode reveals substantially different sink agent behavior (Figure~\ref{fig:runoff_comparison_grid}), particularly under irrigation. Because the MPC controller must predict and avoid waterlogging at sink agents, the cascade mode was adopted to ensure the prediction model accurately reflects within-day runoff accumulation.

\begin{figure}[H]
    \centering
    \includegraphics[width=1\textwidth]{figures/fig_runoff_comparison_grid.png}
    \caption{Hydrological comparison of runoff modes across 0, 5, and 8 mm/day irrigation levels (2024)}
    \label{fig:runoff_comparison_grid}
\end{figure}

\subsection{Scale Justification: Omission of Subsurface Lateral Flow}

The baseline model processes horizontal subsurface lateral flow at the same discrete daily time step as rapid surface runoff. However, given the 460~m$^2$ spatial resolution of the agents (approximately $21 \times 21$~m plots), lateral water movement through dense silty loam over a single day is physically negligible compared to vertical deep drainage and overland flow. 

Therefore, horizontal subsurface flow ($\varphi_4$) was omitted from the proposed framework. The root zones are modeled as independent vertical columns coupled exclusively by the surface runoff routing graph. This physical simplification significantly increases the sparsity of the Jacobian matrices, thereby reducing the computational complexity and accelerating the prediction horizon of the MPC.

\subsection{Agronomic Penalties and Crop Evapotranspiration}

In the original framework, transpiration and drought stress were bounded by the reference evapotranspiration ($\text{ET}_0$), effectively treating all crops equivalently to a standard 12~cm reference grass. This formulation was corrected by introducing crop coefficients ($K_c$) to calculate crop-specific evapotranspiration ($\text{ET}_c$), ensuring Hashemi rice and tobacco respect their distinct biological water demands.

Furthermore, in the original model, the drought stress function $h_3$ appears only in the maturation equation ($x_3$), where it accelerates crop senescence under water deficit. However, $h_3$ is absent from the biomass equation ($x_4$). This means that a completely unirrigated crop would accumulate the same daily biomass increment as a fully irrigated one, provided all other conditions were equal. The only effect of drought would be to advance the onset of senescence, which reduces biomass accumulation only indirectly and with considerable delay.

This omission was identified during the initial ABM validation, when the no-irrigation baseline produced biomass values that were nearly indistinguishable from irrigated scenarios for wheat. The corrected formulation multiplies the radiation use efficiency term ($\theta_{13}$) in the biomass state update by $h_3$, so that drought directly reduces the daily biomass increment. This modification follows the approach of the FAO AquaCrop model, in which water stress reduces radiation use efficiency proportionally to the ratio of actual to potential transpiration. With $\theta_{14} = 0.8$ for rice and 0.6 for tobacco, maximum drought reduces the biomass increment by up to 80\% and 60\%, making irrigation decisions directly consequential for yield outcomes.

This correction is particularly important for the controller comparison objective of this thesis. Without it, the ABM would produce only marginal differences between irrigated and unirrigated outcomes, rendering any comparison between MPC, RL, and baseline controllers meaningless. The original crop used in the baseline paper was wheat in Colombia, which is less water-sensitive (lower $\theta_{14}$); it is possible that the omission went undetected in that context because the irrigation response was already small.

\subsection{Biomass Calibration and Initialization}

The radiation use efficiency parameter $\theta_{13}$ was calibrated so that $x_4$ reliably represents g/m$^2$ of aboveground dry matter. The final agricultural yield is calculated as $Y = x_4 \times \text{HI} \times 10$ (kg/ha). 

For Hashemi rice, the simulated yields were validated against typical Gilan farmer reports ($\sim$2000~kg/ha) and optimal experimental conditions (3500--4000~kg/ha)~\cite{jalali2021}. The initial state conditions reflecting the 20-day nursery period for rice are set to $x_{2,\text{init}} = 210$ and $x_{4,\text{init}} = 60$. Conversely, tobacco seedlings are initialized at $x_{2,\text{init}} = 0$ and $x_{4,\text{init}} = 8$.

\section{Soil and Crop Parameterization}
\label{sec:soil_crop_params}

\subsection{Soil Parameters}
\label{soil_param}
The soil at the study site is classified as silty loam, based on regional soil surveys and FAO soil classification maps for Gilan Province. The soil hydraulic properties used in this study follow FAO Irrigation and Drainage Paper No.~56 recommendations for this soil type. 

\begin{table}[htbp]
\centering
\caption{Soil hydraulic parameters for silty loam (common to all crops)}
\label{tab:soil}

\begin{tabular}{cccc}
\toprule
\textbf{Parameter} & \textbf{Symbol} & \textbf{Value} & \textbf{Source}\\
\midrule
Water uptake coefficient & $\theta_1$ & 0.096 & \cite{Lopezjimenez2024} \\
Wilting point (volumetric fraction) & $\theta_2$ & 0.15 & \cite{allen1998} \\
SCS threshold (mm) & $\theta_3$ & 10.0 & \cite{Lopezjimenez2024} \\
Drainage coefficient & $\theta_4$ & 0.05 & \cite{Lopezjimenez2024} \\
Field capacity (volumetric fraction) & $\theta_6$ & 0.35 & \cite{allen1998} \\
Saturation (volumetric fraction) & $\theta_{sat}$ & 0.50 & \cite{rawls1982} \\
\bottomrule
\end{tabular}
\end{table}

\subsection{Crop Parameters}
\label{crop_param}

Two crops were selected for the simulation environment: rice (Hashemi cultivar) and tobacco (\emph{Nicotiana tabacum}). Both hold significant commercial value in Gilan Province~\cite{surfiran2025, iranica_tobacco}. Their contrasting agronomic characteristics (high water demand, shallow roots, and high drought sensitivity for rice, versus moderate water demand, deep roots, and relative drought tolerance for tobacco) provide a robust basis for evaluating the control algorithms under differing physical constraints.

\subsubsection{Rice (Hashemi Cultivar)}

The base temperature ($T_{\text{base}}$ or $\theta_7$) is set to $10^{\circ}$C, following standard Growing Degree Day (GDD) conventions~\cite{mcmaster1997, laohasiriwong2023} and the baseline ABM~\cite{Lopezjimenez2024}. The maturity threshold ($\theta_{18}$) was calibrated to 1250~$^{\circ}$C$\cdot$day based on the 25-year dataset and validated against local field measurements, which report 1216--1352~GDD for the Hashemi cultivar in Gilan~\cite{sadidi2021}. 

Transplanting is simulated on May 20 (DOY 141), aligning with local practices when the 7-day rolling average $T_{\text{mean}} \geq 18^{\circ}$C and $T_{\min} \geq 12^{\circ}$C~\cite{paredes2025, faghani2011}. The harvest is set for August 20 (DOY 233). Extending the season beyond this date sharply increases the risk of extreme WMO R20mm rainfall events ($\geq$20~mm/day), which can cause crop lodging and severe harvest disruption~\cite{sarma2018, nayak2011}. 

This yields a 93-day field season. As established in Section~\ref{sec:corrections}, the 20-day prior nursery period contributes an initial thermal accumulation of $x_{2,\text{init}} = 210$~GDD and an initial biomass of $x_{4,\text{init}} = 60$~g/m$^2$. Other key parameters include a crop coefficient ($K_c$) of 1.15, a shallow root depth ($\theta_5$) of 400~mm, and a high drought sensitivity ($\theta_{14}$) of 0.80.

\subsubsection{Tobacco (\emph{Nicotiana tabacum})}

The base temperature is similarly set to $10^{\circ}$C, which falls within the standard 10--13$^{\circ}$C restriction range for tobacco~\cite{li2021}. The maturity threshold ($\theta_{18}$) is calibrated to 1200~GDD. This duration corresponds to a 104-day field season, consistent with standard FAO timelines (90--120 days)~\cite{fao_tobacco} and general agronomic data (100--130 days)~\cite{britannica1999}. As tobacco seedlings are typically delivered directly from nurseries with unknown prior conditions, the initial states are set to $x_{2,\text{init}} = 0$ and a minimal $x_{4,\text{init}} = 8$~g/m$^2$.


Transplanting occurs on May 25 (DOY 146), with harvest on September 5 (DOY 249). Tobacco features a significantly deeper root zone ($\theta_5 = 700$~mm) compared to rice, allowing it to access a larger reservoir of deep soil moisture. It also exhibits a lower crop coefficient ($K_c = 0.90$) and lower drought sensitivity ($\theta_{14} = 0.60$), making it more resilient to constrained irrigation schedules. 

\subsubsection{Parameter Summary}

\begin{table}[H]
\centering
\caption{Soil and crop-specific parameters used in the ABM simulation}
\label{tab:params}
\small
\begin{tabular}{llccl}
\toprule
\textbf{Parameter} & \textbf{Symbol} & \textbf{Rice} & \textbf{Tobacco} & \textbf{Source} \\
\midrule
Root depth & $\theta_5$ & 400 mm & 700 mm & \cite{allen1998}, \cite{fao_tobacco} \\
Base temperature & $\theta_7$ & 10$^{\circ}$C & 10$^{\circ}$C & \cite{laohasiriwong2023}, \cite{li2021} \\
Heat stress onset & $\theta_9$ & 35$^{\circ}$C & 30$^{\circ}$C & \cite{fao_tobacco} \\
Extreme heat & $\theta_{10}$ & 42$^{\circ}$C & 38$^{\circ}$C & -- \\
Heat stress maturation rate & $\theta_{11}$ & 0.0030 & 0.0025 & \cite{Lopezjimenez2024} \\
Drought stress maturation rate & $\theta_{12}$ & 0.0030 & 0.0028 & \cite{Lopezjimenez2024} \\
RUE & $\theta_{13}$ & 0.65 & 0.30 & \cite{patel2013}, \cite{ferrara2002} \\
Drought sensitivity & $\theta_{14}$ & 0.80 & 0.60 & \cite{Lopezjimenez2024} \\
Maturity GDD & $\theta_{18}$ & 1250 & 1200 & \cite{sadidi2021}, calib. \\
Max interception & $\theta_{19}$ & 0.95 & 0.90 & \cite{Lopezjimenez2024} \\
50\% intercept. GDD & $\theta_{20}$ & 417 & 450 & \cite{zhang2017} \\
Crop coefficient & $K_c$ & 1.15 & 0.90 & \cite{allen1998}, \cite{fao_tobacco} \\
Depletion fraction & $p$ & 0.20 & 0.50 & \cite{allen1998}, \cite{fao_tobacco} \\
Harvest index & HI & 0.42 & 0.55 & \cite{jalali2021}, \cite{fao_tobacco} \\
Nursery GDD & $x_{2,\text{init}}$ & 210 & 0 & Calib. \\
Initial biomass & $x_{4,\text{init}}$ & 60 g/m$^2$ & 8 g/m$^2$ & Est. \\
\bottomrule
\end{tabular}
\end{table}

\section{Water Budget and Scenario Selection}
\label{sec:water_budget}

\subsection{Water Budget Formulation}

To establish the baseline irrigation requirements for both crops, a seasonal water budget was formulated using the 25-year climatological averages. The total crop water demand ($\text{ET}_c$) is calculated as the product of the average reference evapotranspiration, the crop coefficient, and the season duration:
\begin{equation}
    \text{ET}_c = \overline{\text{ET}_0^{PM}} \cdot K_c \cdot N_{\text{days}}
\end{equation}

The full irrigation requirement ($I_{\text{full}}$) represents the water deficit that must be met by the controller, computed by subtracting the 25-year average seasonal rainfall ($\overline{P}$) from the crop demand:
\begin{equation}
    I_{\text{full}} = \max\!\big(\text{ET}_c - \overline{P},\; 0\big)
\end{equation}

Additionally, constrained irrigation budgets representing 70\% and 50\% of the full requirement were calculated to establish deficit scenarios for evaluating controller robustness under strict water rationing.

\begin{table}[H]
\centering
\caption{Agronomic water budget summary based on 25-year climatology.}
\label{tab:budget}
\begin{tabular}{lcc}
\toprule
\textbf{Parameter} & \textbf{Rice} & \textbf{Tobacco} \\
\midrule
Season duration ($N_{\text{days}}$) & 93 days & 104 days \\
Average ET$_0$ (PM) & 5.02 mm/day & 4.95 mm/day \\
Total crop demand ($\text{ET}_c$) & 537.0 mm & 463.4 mm \\
Average season rainfall ($\overline{P}$) & 53.3 mm & 74.5 mm \\
Full irrigation need ($I_{\text{full}}$) & 483.7 mm & 389.0 mm \\
Deficit budget (70\% of $I_{\text{full}}$) & 338.6 mm & 272.3 mm \\
Deficit budget (50\% of $I_{\text{full}}$) & 241.8 mm & 194.5 mm \\
\bottomrule
\end{tabular}
\end{table}

The calculated rice water demand of 537.0~mm aligns strongly with local Hashemi field measurements, which range from 534~mm under System of Rice Intensification (SRI) management to 632~mm under traditional continuous flooding~\cite{jalali2021}. Similarly, the calculated tobacco demand of 463.4~mm falls perfectly within the expected range of 400--600~mm~\cite{fao_tobacco}. Notably, despite having a longer field season, tobacco requires less total irrigation than rice. This is due to its lower crop coefficient ($K_c$) and the higher average rainfall that occurs during its extended late-summer season.

\subsection{Scenario Selection}

To evaluate controller performance under varying climatic extremes, three specific years were selected from the 25-year dataset to serve as distinct testing scenarios: Dry, Moderate, and Wet. These scenarios are defined by their total seasonal rainfall relative to the historical average.

\begin{table}[H]
\centering
\caption{Selected climate scenarios for controller evaluation.}
\label{tab:scenarios}
\begin{tabular}{lccc}
\toprule
\textbf{Scenario} & \textbf{Year} & \textbf{Rice Rainfall (mm)} & \textbf{Tobacco Rainfall (mm)} \\
\midrule
Dry & 2022 & 39.7 & 54.0 \\
Moderate & 2020 & 42.1 & 73.3 \\
Wet & 2024 & 176.8 & 163.2 \\
\bottomrule
\end{tabular}
\end{table}

\section{Model Validation}
\label{sec:validation}

\subsection{GWETROOT Cross-Validation}

To verify the physical fidelity of the Agent-Based Model's hydrological dynamics, the simulated soil water state ($x_1$) was cross-validated against independent satellite observations. Specifically, the ABM output was compared to the NASA MERRA-2 Root Zone Soil Wetness (GWETROOT) product for the three selected scenario years.

Because the two systems employ different scales and physical boundary conditions, direct absolute comparison is not feasible. Instead, the data were normalized to evaluate the agreement in temporal wetting and drying patterns. The ABM soil water was normalized to a fractional saturation index using the permanent wilting point (WP) and saturation (SAT):
\begin{equation}
    \hat{x}_1(k) = \frac{x_1(k) - \text{WP}}{\text{SAT} - \text{WP}}
\end{equation}
The GWETROOT data is natively provided as a dimensionless 0--1 fraction and was used without modification. Pearson correlation coefficients were computed to assess the temporal alignment between the simulated and observed series.

\begin{table}[H]
\centering
\caption{Cross-validation of ABM soil moisture against NASA GWETROOT.}
\label{tab:crossval}
\begin{tabular}{lccc}
\toprule
\textbf{Scenario} & \textbf{Year} & \textbf{Pearson $r$} & \textbf{$p$-value} \\
\midrule
Dry & 2022 & 0.738 & $3.1 \times 10^{-17}$ \\
Moderate & 2020 & 0.534 & $3.4 \times 10^{-8}$ \\
Wet & 2024 & 0.338 & $9.2 \times 10^{-4}$ \\
\bottomrule
\end{tabular}
\end{table}

\begin{figure}[H]
\centering
\includegraphics[width=\textwidth]{figures/fig_gwetroot_timeseries_pearson.png}
\caption{Normalized soil moisture time series: ABM (blue) vs. NASA GWETROOT (brown).}
\label{fig:gwetroot}
\end{figure}

As shown in Table~\ref{tab:crossval}, the dry-year scenario exhibits a strong positive correlation ($r = 0.738$, $p < 0.001$), confirming that the ABM accurately captures fundamental soil drying and recharge dynamics driven by evapotranspiration and intermittent rainfall. The correlation predictably decreases in the moderate and wet years. This attenuation is expected because frequent rainfall events saturate the soil, reducing the signal contrast and making the correlation metric highly sensitive to minor timing discrepancies in the meteorological data.

While the temporal patterns align, there is a consistent offset in the absolute levels; the normalized ABM baseline fluctuates near 0.07, whereas the GWETROOT average sits higher at approximately 0.47. This discrepancy is attributable to three distinct structural differences between the datasets:
\begin{enumerate}
    \item \textbf{Observation Depth:} NASA MERRA-2 models a deep root zone of roughly 1000~mm, which retains moisture significantly longer than the ABM's shallow 400~mm rice root zone.
    \item \textbf{Irrigation Status:} The GWETROOT satellite product observes actual field conditions in Gilan, which includes widespread anthropogenic irrigation. Conversely, this specific ABM validation run simulates an unirrigated, rain-fed control scenario to establish a zero-control baseline.
    \item \textbf{Soil Parameterization:} The standardized MERRA-2 global soil grid relies on generalized hydraulic parameters that differ from the highly specific silty loam parameters calibrated for this local model.
\end{enumerate}

\section{Conclusion}
\label{sec:conclusion}

This chapter established a high-fidelity, mathematically rigorous virtual environment for simulating agricultural dynamics in Gilan Province. By building upon an existing topographical agent-based framework, several critical extensions were introduced to ensure the model's suitability for advanced control engineering applications. 

Most notably, the hydrological dynamics were decoupled into distinct surface and subsurface layers. This architectural correction eliminated mass-balance violations present in the baseline literature, ensuring strict conservation of water fluxes and physically realistic surface pooling. Furthermore, the direct integration of drought stress penalties into the biomass accumulation equation established a necessary, responsive cost function, ensuring that simulated crop yields dynamically reflect the success or failure of applied irrigation strategies.

The environment was successfully parameterized for two commercially significant local crops, the highly water, sensitive Hashemi rice and the more resilient tobacco, using a 25-year localized climatology. Finally, the baseline temporal dynamics of the model were successfully cross-validated against independent NASA MERRA-2 satellite observations.

With the physical, biological, and mathematical boundaries of the virtual plant now validated, this environment satisfies its dual objective. It serves both as the dynamic plant for the Model Predictive Controller and as the episodic training environment for the Reinforcement Learning agent. The subsequent works will detail the formulation, training, and comparative evaluation of these control algorithms across the established dry, moderate, and wet climatic scenarios.